\documentclass[11pt]{article}
\usepackage[margin=1in]{geometry}
\usepackage[T1]{fontenc}
\usepackage[utf8]{inputenc}
\usepackage{amsmath,amssymb,amsthm,mathtools}
\usepackage{enumitem}
\usepackage{hyperref}
\usepackage{listings}
\lstset{basicstyle=\ttfamily\small,columns=fullflexible,breaklines=true,
  literate={∀}{{$\forall$}}1 {∃}{{$\exists$}}1 {→}{{$\to$}}1 {¬}{{$\neg$}}1
           {≤}{{$\leq$}}1 {≥}{{$\geq$}}1 {∧}{{$\wedge$}}1 {∨}{{$\vee$}}1
           {ℚ}{{$\mathbb{Q}$}}1 {ℤ}{{$\mathbb{Z}$}}1 {∣}{{$\mid$}}1}

\newtheorem{theorem}{Theorem}
\newtheorem{proposition}{Proposition}
\newtheorem{lemma}{Lemma}
\newtheorem{definition}{Definition}
\newtheorem{remark}{Remark}

\title{A Lean-oriented core for the $Q_{a,b}$ package-coprimality proof\\
       Version 4: worker handoff and polished assumption boundary}
\author{Draft blueprint for review}
\date{\today}

\begin{document}
\maketitle

\begin{abstract}
This note extracts a Lean-oriented core from the current proof of the
$Q_{a,b}$ package-coprimality theorem.  It is not intended to replace the full
manuscript.  Its purpose is to define a conditional formalization target with
explicit theorem packs that are generous enough for a first Lean
implementation, but narrow enough that later phases can open the packs and
formalize their contents down to Mathlib and Lean-checkable certificates.

Version 4 incorporates the final polish review of the v3 package.  The residual
witness name is made Lean-friendly, the certificate interfaces no longer import
the core proof, the lower-box hypothesis is passed explicitly to the lower
finite eliminators, and the phase-1 polynomial semantics now distinguish
primitive orientations, full imprimitive orientations, reciprocal package
products, and orientation-level versus product-level sharing.
\end{abstract}

\section{Purpose and scope}

The full manuscript contains discovery scaffolding and auxiliary finiteness
arguments that are not needed in the shortest Lean validation path.  The core
retained here is the following:
\[
\begin{gathered}
\text{package sharing}
\Rightarrow \text{normalized counterexample},\\
D\leq 175{,}394{,}637,\\
D\geq 6{,}816{,}242 \Rightarrow \text{upper unit-normalized contradiction},\\
D\leq 6{,}816{,}241 \Rightarrow \text{admissible lower state},\\
\text{unit/both-nonunit/one-nonunit branch eliminations},\\
\text{residual one-nonunit envelope and terminal certificates}.
\end{gathered}
\]
The top-level Lean theorem will be conditional on named theorem packs.  The
computational eliminations should ultimately be discharged by Lean-checkable
certificates rather than by trusting external loops.

The following ingredients are deliberately absent from the top-level core:
ASZ-style unlikely-intersection boundedness, Beukers--Schlickewei S-unit
finiteness, density-one isolation side results, and most sign/phase/Schur
scaffolding.  These may remain in the human manuscript as explanation or
independent checks, but they are not part of the shortest Lean proof spine.

\section{Constants and normalized counterexamples}

We write
\[
B_0=175{,}394{,}637,
\qquad
B_1=6{,}816{,}241,
\qquad
B_1^+=6{,}816{,}242,
\]
\[
R=16{,}583,
\qquad
R^+=16{,}584,
\qquad
E_R=2601.
\]
Here $B_0$ is the global height--degree--Mahler bound, $B_1^+$ is the first
upper-range value, $B_1$ is the lower-box cutoff, $R^+$ is the start of the
large one-nonunit branch, $R$ is the residual cutoff, and $E_R$ is the residual
degree bound.

\begin{definition}[Normalized counterexample record]
A Lean record \verb|NormCE| contains natural numbers
\[
 m<n,
 \qquad p<q,
 \qquad D,d,U,V_0,
\]
with $m,p,d,U,V_0>0$, together with
\[
D=\max\{m,n,p,q\},
\qquad
\gcd(\gcd(m,n),\gcd(p,q))=1,
\]
and a proof that the two normalized shapes are not the same up to swapping.
Here $d$ is the degree of the selected primitive minimal polynomial, and
$U,V_0$ are the absolute endpoint coefficients used in the finite branch
analysis.
\end{definition}

Define
\[
\begin{aligned}
\operatorname{Unit}(ce)&:\Longleftrightarrow U=1\text{ and }V_0=1,\\
\operatorname{BothNonunit}(ce)&:\Longleftrightarrow 2\leq U\text{ and }2\leq V_0,\\
\operatorname{OneNonunit}(ce)&:\Longleftrightarrow
   (U=1\text{ and }2\leq V_0)\text{ or }(2\leq U\text{ and }V_0=1).
\end{aligned}
\]
Because $U,V_0>0$, the trichotomy
\[
\operatorname{Unit}(ce)\vee \operatorname{BothNonunit}(ce)\vee
\operatorname{OneNonunit}(ce)
\]
is elementary arithmetic.  It is proved directly in Lean and is not part of the
local p-adic theorem pack.

\section{Opaque selected counterexamples}

The selected counterexample predicate is
\[
\texttt{FromShare}(P,Q,ce).
\]
For the first conditional formalization it is an opaque predicate.  It must not
be defined as \verb|True|.  Otherwise broad finite-elimination assumptions would
apply to arbitrary dummy records and could make the development inconsistent.

In a later phase, \verb|FromShare| should become a structure containing the
actual selected orientations, an off-unit-circle common root, the primitive
minimal polynomial, degree and endpoint-coefficient data, and divisibility by
the selected package polynomials.

\section{Package polynomial semantics}

The Lean definition of \verb|PackageShare| is now product-level sharing of the
reciprocal package products.  This replacement is deliberately careful about
primitive versus imprimitive pairs.  The coefficient formula below is for the
primitive orientation polynomial only.  For coprime positive $A,B$, define
\[
Q^{\rm prim}_{A,B}(x)=
\sum_{0\leq j<A} B(j+1)x^j+
\sum_{A\leq j\leq A+B-2} A(A+B-j-1)x^j.
\]
For a nonprimitive pair $a=gA$, $b=gB$, with $g=\gcd(a,b)$, the full orientation
polynomial should use the manuscript convention
\[
Q_{a,b}(x)=g Q^{\rm prim}_{A,B}(x^g).
\]
The reciprocal package product is
\[
P_{a,b}(x)=Q_{a,b}(x)Q_{b,a}(x).
\]
Thus the Lean development should keep separate names:
\[
\texttt{qPrimZ},\quad
\texttt{qOrientZ},\quad
\texttt{qPackageProdZ},\quad
\texttt{OrientationShare},\quad
\texttt{PackageShare}.
\]
Product-level \verb|PackageShare| is closest to the theorem statement.  If the
collision dictionary uses orientation-level sharing, it should prove a separate
bridge from product sharing to an orientation share for some reciprocal-choice
of orientations.  The current Lean API proves this bridge by passing to an
irreducible factor of the common product-level witness.

\section{The theorem packs}

\subsection{Collision pack}

\begin{lstlisting}[language={}]
structure CollisionPack where
  share_to_normCE :
    ∀ {P Q : PosPair},
      P.unequal → Q.unequal → ¬ P.sameUnordered Q →
      PackageShare P Q → ∃ ce : NormCE, FromShare P Q ce
\end{lstlisting}

This pack should eventually be opened by formalizing the package polynomial
layer, product-to-orientation sharing, the collision dictionary, removal of
forced cyclotomic roots, reciprocal orientation handling, no-three
compatibility, and primitive normalization.

\subsection{Global height pack}

\begin{lstlisting}[language={}]
structure GlobalHeightPack where
  absolute_bound :
    ∀ {P Q : PosPair} {ce : NormCE},
      FromShare P Q ce → ce.D ≤ 175394637
\end{lstlisting}
It packages the height of a collision root, the rank-two torus/isolation degree
bound, the Mahler/Smyth gap or replacement, and the exact constant arithmetic.
It is the part of the proof that makes the infinite family finite.

\subsection{Upper normalization and upper elimination}

\begin{lstlisting}[language={}]
structure UpperNormalizationPack where
  upper_unit_normalized :
    ∀ {P Q : PosPair} {ce : NormCE},
      FromShare P Q ce → 6816242 ≤ ce.D → ce.UnitCoeff

structure UpperRangePack where
  no_upper_unit_counterexample :
    ∀ {P Q : PosPair} {ce : NormCE},
      FromShare P Q ce → ce.UnitCoeff →
      6816242 ≤ ce.D → ce.D ≤ 175394637 → False
\end{lstlisting}
This split reflects the proof more faithfully than a single broad upper-range
contradiction.  Later, the second pack can be replaced by coverage certificates
and terminal row exclusions.

\subsection{Local lower-state extraction}

\begin{lstlisting}[language={}]
structure LowerState where
  rowId : Nat

constant Encodes : NormCE → LowerState → Prop
constant AdmissibleLowerState : LowerState → Prop

structure LocalPacketPack where
  to_lower_state :
    ∀ {P Q : PosPair} {ce : NormCE},
      FromShare P Q ce → ce.D ≤ 6816241 →
      ∃ st : LowerState, Encodes ce st ∧ AdmissibleLowerState st
\end{lstlisting}
In later phases, \verb|LowerState| should expose primitive shapes, scales,
extraction degrees, endpoint coefficients, slope signatures, packet
occupancies, deficient-prime data, total-radical constraints, and row IDs.

\subsection{Residual envelope}

For the residual one-nonunit branch, the local pack also exposes reciprocal
normalization and the exact residual envelope:
\[
D\leq16{,}583,
\qquad
V_0=1,
\qquad
2\leq U\leq D,
\qquad
d\leq2601.
\]
In Lean:
\begin{lstlisting}[language={}]
def ResidualEnvelope (ce : NormCE) : Prop :=
  ce.D ≤ 16583 ∧
  ce.Vabs = 1 ∧
  2 ≤ ce.U ∧ ce.U ≤ ce.D ∧
  ce.d ≤ 2601

structure ResidualWitness (ce : NormCE) where
  ce' : NormCE
  st : ResidualState
  recnorm : ReciprocalNormalize ce ce'
  envelope : ResidualEnvelope ce'
  encodes : EncodesResidual ce' st
  admissible : AdmissibleResidualState st
\end{lstlisting}
This prevents the residual eliminator from hiding the reductions to $|V|=1$,
$2\leq U\leq D$, and $d\leq2601$.

\subsection{Finite lower eliminations}

\begin{lstlisting}[language={}]
structure FiniteLowerPack where
  no_lower_unit :
    ∀ {ce : NormCE} {st : LowerState},
      Encodes ce st → AdmissibleLowerState st →
      ce.D ≤ lowerMax → ce.UnitCoeff → False

  no_both_nonunit :
    ∀ {ce : NormCE} {st : LowerState},
      Encodes ce st → AdmissibleLowerState st →
      ce.D ≤ lowerMax → ce.BothNonunit → False

  no_one_nonunit_large :
    ∀ {ce : NormCE} {st : LowerState},
      Encodes ce st → AdmissibleLowerState st → ce.OneNonunit →
      16584 ≤ ce.D → ce.D ≤ 6816241 → False
\end{lstlisting}
The explicit lower-box argument in the first two fields is deliberate.  The
state should also encode the lower box eventually, but the core proof already
has the bound and passes it to make the finite-pack boundary auditable.

\subsection{Residual finite elimination}

\begin{lstlisting}[language={}]
structure FiniteResidualPack where
  no_residual_state :
    ∀ {ce ce' : NormCE} {st : ResidualState},
      ReciprocalNormalize ce ce' →
      ResidualEnvelope ce' →
      EncodesResidual ce' st →
      AdmissibleResidualState st →
      False
\end{lstlisting}
The first concrete certificate milestone should target this pack.

\section{Core conditional theorem}

\begin{theorem}[Core assembly theorem]
Assume the seven theorem packs above.  Then package sharing cannot occur
between two distinct nontrivial reciprocal packages.
\end{theorem}

\begin{proof}
Assume a package share between two distinct nontrivial packages.  The collision
pack gives a normalized counterexample $ce$ with \verb|FromShare P Q ce|.  The
global height pack gives $D\leq B_0$.

If $D\geq B_1^+$, upper normalization gives unit endpoint coefficients, and the
upper-range pack gives a contradiction.

Thus $D\leq B_1$.  The local packet pack gives an admissible lower state
encoding $ce$.  By the arithmetic endpoint trichotomy, $ce$ is either unit,
both-nonunit, or one-nonunit.  The first two cases are eliminated by the lower
finite pack using the explicit hypothesis $D\leq B_1$.

In the one-nonunit case, if $D\geq R^+$, the large one-nonunit eliminator gives
a contradiction.  Otherwise $D\leq R$.  The local residual theorem supplies a
reciprocally normalized residual witness satisfying
$D\leq R$, $V_0=1$, $2\leq U\leq D$, and $d\leq E_R$.  The residual finite pack
eliminates this state.  Hence no normalized counterexample exists, and the
original package share was impossible.
\end{proof}

\section{Certificate strategy}

The computational part should ultimately be Lean-checked by certificates, not
trusted external loops.  The certificate order should be:
\begin{enumerate}[label=(\arabic*)]
\item implement canonical primitive and imprimitive package polynomial
      constructors;
\item formalize a good-reduction theorem for primitive common factors;
\item check the residual terminal modular gcd certificates over
      $\mathbf F_{1009}$;
\item add residual coverage certificates for the state generator and pair
      sieve;
\item extend the same pattern to the both-nonunit, one-nonunit-large, and unit
      branches.
\end{enumerate}

The schematic certificate interface now includes a terminal-row well-formedness
predicate, an \verb|exactGcdDegree| field, and a tag specifying whether the
certificate checks divided package polynomials or undivided collision
trinomials.  For modular gcd certificates, the proof must specify which family
is checked.  If the certificate reports modular gcd degree $2$ for collision
trinomials, the degree $2$ must be explicitly accounted for by the forced
double root at $x=1$.

\section{Implementation phases}

\begin{enumerate}[label=\textbf{Phase \arabic*.},leftmargin=*]
\item \textbf{Compile the conditional scaffold.}  Build the Lake project, fix
minor API drift, and keep the core theorem pack-parametric.

\item \textbf{Define package polynomials and package sharing.}  Implement the
primitive coefficient polynomial, the full imprimitive orientation polynomial,
the reciprocal product, and product-level \verb|PackageShare|.

\item \textbf{Open the residual terminal certificates.}  Encode the small
terminal residual rows and the modular gcd certificates over $\mathbf F_{1009}$.

\item \textbf{Open residual coverage.}  Expand \verb|ResidualState|, define row
predicates, and verify coverage of the residual generator and pair sieve.

\item \textbf{Open lower nonunit branches.}  Formalize the both-nonunit and
large one-nonunit finite-state eliminations.

\item \textbf{Open the unit branch.}  Formalize the full-Kummer terminal
certificates and decide whether the Kummer-defect branch remains a narrow
isolation axiom or is replaced by direct finite certificates.

\item \textbf{Open the collision pack.}  Replace \verb|FromShare| by real
selected-root data and prove collision extraction.

\item \textbf{Open the height and local theorem packs.}  Formalize or narrowly
axiomatize the height--degree--Mahler argument, Borisov/Newton consequences,
and Kummer extraction-degree control.
\end{enumerate}

\section{Open mathematical review questions}

\begin{enumerate}[label=(\arabic*)]
\item Should the unit Kummer-defect branch keep a narrow global-isolation axiom
for the 46 irreducible shapes, or should it be replaced by direct finite
certificates?
\item Should terminal modular gcd certificates be attached to divided package
polynomials or to undivided collision trinomials?
\item Should the upper-range final rows be excluded by the same-shape degree
contradiction or by terminal modular certificates?
\item What is the best eventual representation of \verb|FromShare|: explicit
field/root data, minimal-polynomial data, or divisibility-only data?
\end{enumerate}

\section{Acceptance criterion for the first Lean milestone}

The first milestone is complete when:
\begin{enumerate}[label=(\arabic*)]
\item \verb|lake build| succeeds;
\item \verb|package_coprimality_from_packs| compiles;
\item \verb|FromShare| is opaque or a genuine structure, not \verb|True|;
\item all broad temporary assumptions live only in \verb|BroadAxioms.lean|;
\item \verb|Qab.lean| does not import \verb|BroadAxioms.lean|;
\item the README explicitly says the project is conditional on theorem packs;
\item the core proof file remains a short assembly theorem.
\end{enumerate}

\end{document}
